\section{Closed routes, with the mechanism that closed them}

This section catalogues every no-go, countermodel, obstruction, and refuted
proof strategy found in the \#249/\#257 corpus that bears on Erd\H{o}s \#249
(\(S := \sum_{n\ge1}\varphi(n)/2^n\) irrational). Erd\H{o}s \#249 is
\ev{Open}; nothing in this section decides it. What follows is a machinery
inventory: for each dead route, the exact statement that closed it, the exact
scope of the closure, and the exact machinery salvaged from the wreckage.
The organising warning, stated once here and applicable to every item below:
several of these ``obstructions'' are statements about a
\emph{representation} of the problem (a proof strategy, a certificate
family, a coordinate system) and not about the object \(S\) itself. Confusing
the two is the single most common error this catalogue exists to prevent.

\subsection{Meta-level: finite inspection cannot certify the supply (the
\texorpdfstring{\(\gamma\)}{gamma}-splice remark)}

\textbf{(a) Route as conceived.} Every unconditional route to \#249 in this
corpus bottoms out in a \emph{certificate-supply} obligation of the shape
\(\forall h\ge1\,\forall N_0\,\exists N\ge N_0\,\exists L,\ \mathrm{Sep}(h,N,L)\)
(\lean{Erdos249257.irrational\_totient\_series\_iff\_certificate\_supply}{LcmConeFlatness.lean:412}, the exact
open-obligation form quoted throughout the existing manuscript). A natural
hope is that verifying \(\mathrm{Sep}\) — or its diagonal specialisation — at
every scale up to some large, explicit bound \(B\) is evidence that the
supply holds, and that pushing \(B\) far enough would eventually amount to a
proof.

\textbf{(b) Exact mechanism that closed it.} The manuscript's own remark
(\ev{Math}, prose-only, not formalised in Lean:
\emph{\S5.4, ``finite inspection cannot establish the supply''}) gives an
explicit splice construction. Take any coefficient stream \(\gamma:\mathbb
N\to\mathbb N\) that agrees with \(\varphi\) on every index \(\le B\), and
alter \(\gamma\) at one residue class beyond \(B\) so that the resulting
binary series \(\sum\gamma(n)/2^n\) is forced \emph{rational}. Such a
\(\gamma\) exists for every \(B\) (this is the same coboundary-splice
technique that produces \S\ref{sec:parity-countermodel}'s
\texttt{parityCoboundaryWeight} witness, generalised: insert a lacunary
zero-valued coboundary edit \(2/2^m - 4/2^{m+1} = 0\) at some \(m>B\)). Because
\(\gamma\) agrees with \(\varphi\) on every index \(\le B\), \(\gamma\) passes
every finite certificate that only inspects indices \(\le B\) — in
particular every instance of \(\mathrm{Sep}(h,N,L)\) with \(N+L\le B\) that
\(\varphi\) itself passes or fails. Yet \(\sum\gamma(n)/2^n\in\mathbb Q\) by
construction.

\textbf{(c) Precise scope.} This is a statement about \emph{proof method},
not about \(\varphi\): it does not touch \(S\) at all. What it excludes is
the inference ``\(\mathrm{Sep}\) verified up to bound \(B\), for arbitrarily
large \(B\), therefore \(\mathrm{Sep}\) holds cofinally.'' The scope is
exactly \(\scale{bounded}\) versus \(\scale{cofinal}\): any finite-B
verification, however large, is compatible with both outcomes, because a
rational stream can be built to survive any \emph{fixed} inspection horizon.
It says nothing about whether \(\varphi\) itself is such a \(\gamma\) — only
that no finite-horizon check can distinguish \(\varphi\) from a \(\gamma\)
that is.

\textbf{(d) Salvage.} The splice construction is exactly the mechanism
underlying \S\ref{sec:parity-countermodel}'s Lean-formalised countermodel
below, so its \emph{content} is not lost even though this specific remark is
unformalised: everywhere a proof strategy in this corpus proposes to certify
a supply by exhaustive finite search (the diagonal deposits at
\(t\in\{1,\dots,64\}\), \lean{Erdos249257.certifiedKill\_diagonal\_all\_imported\_through\_t64}{DiagonalPincerCertificatesT64.lean:1967}; the row-200{,}000 sqrt-escape
census on the \#257 side), this remark is the standing reason such a search,
however large, is evidence and not proof. \ev{Math} \scale{n/a}
\coord{binary-digit}.

\subsection{The parity-aperiodicity countermodel
(\texttt{TotientParityCoboundaryCountermodel})}
\label{sec:parity-countermodel}

\textbf{(a) Route as conceived.} A natural strengthening of \#249 attempts:
show that any coefficient word \(c:\mathbb N\to\mathbb N\) that (i) is
uniformly bounded, (ii) satisfies the trivial linear-growth bound
\(c(n)\le n\), (iii) matches \(\varphi(n)\bmod 2\) exactly, and (iv) is not
eventually periodic — even strengthened to \emph{cofinally, arbitrarily
separated, arbitrarily long blocks} of non-periodicity — must have irrational
binary series \(\sum c(n)/2^n\). This is the natural target for anyone
trying to use only \(\varphi\)'s coarse parity/aperiodicity profile.

\textbf{(b) Exact mechanism.} The construction
\(\mathtt{parityCoboundaryWeight}(n) := \mathtt{parityBaseWeight}(n) +
2\cdot\mathtt{largePowerTwoBit}(n) - 4\cdot\mathtt{largePowerTwoBit}(n-1)\),
where \(\mathtt{parityBaseWeight}\) is the eventually-constant word
\(0,1,1,2,4,4,4,\dots\) and \(\mathtt{largePowerTwoBit}(n)=1\) iff
\(n=2^{k+3}\), is a lacunary zero-valued coboundary splice on top of an
eventually-constant base. The strongest of five landed theorems,
\begin{prop}[Totient-parity arbitrarily-many-separated-carry rational
countermodel]
There exists \(c:\mathbb N\to\mathbb N\) such that: \(c(n)\le6\) for all
\(n\); \(c(n)\le n\); \(c(n)\equiv\varphi(n)\pmod2\) for every \(n\); for
every \(N,G,K\) there is a block of \(K\) explicit \((6,0)\) carry-pulse
pairs beyond \(N\), each pair separated by more than \(G\); and
\(\sum_n c(n)/2^n = 3/2\).
\end{prop}
\lean{Erdos257PeriodNoncollapse.exists\_totientParity\_arbitrarilyManySeparatedCarry\_rational\_countermodel}{TotientParityCoboundaryCountermodel.lean:637}
(sum computation \lean{tsum\_parityCoboundaryWeight\_eq\_three\_halves}{TotientParityCoboundaryCountermodel.lean:359}; non-periodicity
\lean{parityCoboundaryWeight\_not\_eventually\_periodic}{TotientParityCoboundaryCountermodel.lean:487}; exact parity match
\lean{parityCoboundaryWeight\_mod\_two\_eq\_totient}{TotientParityCoboundaryCountermodel.lean:194}). \ev{Lean} \scale{cofinal}
\coord{binary-digit}.

\textbf{(c) Precise scope.} This is an \emph{object}-level existence witness
(\(c\ne\varphi\), only \(c\equiv\varphi\pmod2\)) that closes a
\emph{proof-route}: it proves that the hypothesis set \(\{\)bounded, linear
growth, \(\varphi\)-parity match, non-eventual-periodicity\(\}\), no matter
how strongly the non-periodicity clause is strengthened (up to the
arbitrarily-separated arbitrarily-numerous form above), can never entail
irrationality of the associated binary series, because \(c\) satisfies every
hypothesis and is rational. It says nothing about \(\varphi\) itself — the
witness sequence \(c\) is a different, hand-built sequence. Any future
sufficient-condition candidate stated purely in terms of \(\{\)coefficient
boundedness, growth, parity, periodicity\(\}\) must be checked against this
countermodel before being trusted, for \emph{either} \#249 or \#257.

\textbf{(d) Salvage.} The lacunary coboundary-splice technique itself —
inserting zero-valued edits \(2\cdot2^{-m}-4\cdot2^{-(m+1)}=0\) at sparse
ranks to destroy periodicity while preserving the rational sum — is a fully
general recipe for building rational-valued, non-eventually-periodic,
bounded coefficient sequences matching \emph{any} parity template, and is
literally the mechanism instantiated informally in
Theorem~\ref{thm:gamma}'s \(\gamma\)-splice construction. It is also a
ready-made stress-test fixture for any future parity-based sufficient
condition proposed anywhere in the corpus. Confirms the corpus-wide lesson
(project memory \texttt{feedback\_erdos\_reductions\_rejected\_bank\_real\_results}):
only arguments using \emph{actual} quantitative totient/Mersenne size or
residue information — as in the \texttt{TotientActualLcm*},
\texttt{TotientFixedRank*} families below — can possibly close \#249; pure
symbolic-word arguments cannot.

\subsection{Two scoped \#249 no-go countermodels: square-CRT correction
suppression is ambiguous}

\textbf{(a) Route as conceived.} The totient prime-dilation cocycle
\(\varphi(pm)=(p-1)\varphi(m)+[p\mid m]\varphi(m)\) carries a correction term
whenever \(p\mid m\). A natural strategy imposes a square congruence
\(n\equiv A+pa\pmod{p^2}\) with a ``clean'' shift \(h\) (\(p\nmid a+h\)) to
suppress that correction on a finite horizon, then hopes that
correction-suppression (``cleanliness'') is itself enough to guarantee that
the resulting finite-difference cube is \emph{nonzero} — i.e. that cleaning
away the correction term automatically leaves behind a genuine, nonvanishing
arithmetic signal usable in a curvature argument.

\textbf{(b) Exact mechanism that closed it.} Two explicit, opposite
finite computations refute the hoped-for entailment simultaneously:
\begin{itemize}[nosep]
\item \emph{Vanishing witness.} The smallest returned countermodel has the
\emph{entire} two-step finite block \(\mathtt{actualOneCubeCoeff}\) equal to
zero at \(n=52\), \(r_0=13\), \(r_1=13/18\)
(\lean{squareCRT\_clean\_block\_can\_vanish}{SquareCRTCube.lean:459}, kernel-\texttt{decide}d;
cleanliness witnessed by \lean{squareCRT\_vanishing\_countermodel\_is\_clean}{SquareCRTCube.lean:468}).
\item \emph{Nonvanishing witness.} A separate, nearby clean configuration at
\(n=27\) has the same block \emph{nonzero}
(\lean{squareCRT\_clean\_block\_can\_be\_nonzero}{SquareCRTCube.lean:477}; cleanliness witnessed by
\lean{squareCRT\_nonzero\_countermodel\_is\_clean}{SquareCRTCube.lean:485}).
\end{itemize}
Both witnesses are exhibited alongside the positive achievability result
\lean{exists\_squareCRT\_clean\_horizon}{SquareCRTCube.lean:327} (a bounded common base satisfying prescribed
CRT anchors/residues does exist). \ev{Cert} (theorem bodies not
independently re-verified beyond signatures and docstring — flagged
\texttt{DOCSTRING-SOURCED} in the source bank) \scale{fixed}
\coord{other:square-crt-cocycle}.

\textbf{(c) Precise scope.} These are two \#249-scoped countermodels — the
coordinate is the totient prime-dilation cocycle specifically, not a
problem-agnostic phenomenon. Together they show that
\emph{correction-suppression is logically independent of the
sign/nonvanishing question it was hoped to settle}: ``clean'' (i.e.
correction-free) is consistent with \emph{both} outcomes. This is a
statement about a specific finite-difference construction (the
representation), not about \(\varphi\): it does not exclude a nonvanishing
finite-difference cube in general, only the inference from cleanliness
\emph{alone} to nonvanishing.

\textbf{(d) Salvage.} The general cocycle identity
\(\varphi(pm)=(p-1)\varphi(m)+[p\mid m]\varphi(m)\)
(\lean{totient\_mul\_prime\_cocycle}{SquareCRTCube.lean:38}) and the CRT-witness apparatus
(\lean{exists\_squareCRT\_base}{SquareCRTCube.lean:230}, \lean{exists\_large\_squareCRT\_base}{SquareCRTCube.lean:263}) remain fully reusable for
\emph{any} residue-forcing construction, in any coordinate. More valuably,
the abstract cube-cancellation lemma extracted from the same file,
\lean{alternating\_powerset\_sum\_eq\_zero\_of\_insert\_invariant}{SquareCRTCube.lean:363} — ``if inserting one
coordinate never changes the value of \(f\), the alternating powerset sum of
\(f\) vanishes'' — is completely divorced from totient/CRT content and is
directly reusable for any Hilbert-cube or higher-difference argument on
either problem. What remains as the honest open gap: an independent
anti-concentration or residue-producer argument for why the clean residual
does \emph{not} vanish, which this module explicitly does not supply.

\subsection{The adelic height obstruction}
\label{sec:adelic-height}

\textbf{(a) Route as conceived.} Several proof strategies attempt to
\emph{clear} an unwanted denominator factor from a rational value \(x\) by
multiplying by a small integer coefficient \(c\), hoping to discard the
complementary denominator information entirely (e.g. localising a diagonal
tail difference to one LCM ray channel while dropping the rest).

\textbf{(b) Exact mechanism.}
\begin{lem}[Scalar-localisation complement divisibility]
For \(x:\mathbb Q\), \(c:\mathbb Z\), \(H:\mathbb N\): if \(H\mid x.\mathrm{den}\)
and \((c\cdot x).\mathrm{den}\mid H\), then \(x.\mathrm{den}/H \mid
c.\mathrm{natAbs}\).
\end{lem}
\lean{Erdos249257.AdelicHeightObstruction.scalarLocalization\_complement\_dvd}{AdelicHeightObstruction.lean:23}. The equality
strengthening, \(\exists t:\mathbb Z,\ H\cdot c\cdot x = t\cdot x.\mathrm{num}\), is
\lean{Erdos249257.AdelicHeightObstruction.scalarLocalization\_integer\_eq\_mul\_num}{AdelicHeightObstruction.lean:56}. The
Mersenne-scale corollary — for positive \(x\), if \(2^r\mid x.\mathrm{num.natAbs}\)
and \(x<2/(2^n-1)\), then \(2^r(2^n-1)<2\cdot x.\mathrm{den}\) — is
\lean{Erdos249257.AdelicHeightObstruction.positiveRat\_mersenne\_height}{AdelicHeightObstruction.lean:103} (generic form
\lean{positiveRat\_numDivisor\_mul\_lt\_two\_mul\_den}{AdelicHeightObstruction.lean:77}).

A companion rigidity fact
rules out smuggling the discarded information back in through a linear
channel: if a linear map \(\Lambda:V\to W\) factors through a surjective
evaluation \(\mathrm{ev}\) (\(\ker\mathrm{ev}\le\ker\Lambda\)), then
\(\Lambda\) is \emph{exactly} \(\mathrm{ev}(\cdot)\bullet w_0\) for a single
fixed \(w_0\) —
\lean{Erdos249257.AdelicHeightObstruction.linearDescender\_eq\_smul\_eval}{AdelicHeightObstruction.lean:120}.

The upstream
primitive these import,
\lean{Erdos249257.RationalDenominatorSurvival.divisor\_dvd\_divInt\_den}{RationalDenominatorSurvival.lean:17} and
\lean{survivingDivisor\_dvd\_scaled\_divInt\_den}{RationalDenominatorSurvival.lean:38}, gives the exact
survival/cancellation law: a divisor \(m\mid D\) of a displayed denominator
survives reduction of \(a/D\) iff \(m\) is coprime to \(a\); after scaling the
numerator by \(h\), a surviving coprime divisor \(C\mid D\) shrinks to
exactly \(C/\gcd(C,h)\), never further. \ev{Lean} \scale{uniform}
\coord{other:rational-height}.

\textbf{(c) Precise scope.} This is pure \(\mathbb Q\)-arithmetic with zero
Mersenne or totient content — it is not a statement about \(S\) at all, but
about \emph{any} rational-height bookkeeping argument. Its exact content:
scalar denominator-clearing never \emph{erases} the complementary
denominator; it moves that complement into the coefficient's (Archimedean)
size. Any construction that tries to shrink a displayed denominator down to
one surviving channel pays for it in coefficient growth. This directly
explains \emph{why} denominator-compression strategies are structurally hard:
the Mersenne-shadow denominator lower bound
\(2^{t/2}\le\prod_{p\in\mathrm{upperHalfPrimes}\,t}\mathrm{mersenne}(p)\)
(\lean{upperHalfMersenneProduct\_lower\_bound}{MersenneShadowDenominatorGrowth.lean:60}) forces a
quadratic-in-scale coefficient cost via this lemma if one tries to localise
to a single channel.

\textbf{(d) Salvage.} Fully problem-agnostic and directly reusable for \#257
or any other Lambert-type denominator-survival argument as-is. It is the
general mechanism underlying the corpus's Farey-gap denominator-exclusion
family (Part C of the certificate bank,
\lean{tsum\_totient\_div\_pow\_two\_ne\_ratCast\_of\_den\_le\_79639646646701375323355774875831053}{CertificateKernel.lean:18384}, the strongest
current unconditional statement about \(S\): if \(S\) is rational its
reduced denominator exceeds \(79{,}639{,}646{,}646{,}701{,}375{,}323{,}355{,}774{,}875{,}831{,}053\approx7.96\times10^{34}\)).

\subsection{The signed \(q\)-moment machinery — a positive tool, not an
obstruction}

\textbf{(a) Route as conceived.} A ``signed Hankel determinant'' /
Hermite--Pad\'e style route to \#249 would exhibit a finite signed linear
combination of dyadically-cleared totient terms and show it is provably
nonzero, giving a nonvanishing certificate for a first-harmonic or
Möbius-companion construction.

\textbf{(b) Exact mechanism (a nonvanishing \emph{engine}, not a no-go).}
\begin{lem}[Unique-terminal parity nonvanishing]
If a finite signed sum \(\sum_{i\in s}u(i)\cdot2^{e(m)-e(i)}\), clearing
dyadic denominators to a common exponent \(e(m)\), has one index \(m\in s\)
with strictly maximal exponent \(e(m)\) and odd coefficient \(u(m)\), while
every other index has strictly smaller exponent, then the whole cleared sum
is \(\equiv1\pmod2\), hence nonzero.
\end{lem}
\lean{scaled\_dyadic\_sum\_ne\_zero}{SignedQMomentObstruction.lean:96}, built from a fully generic rectangular
Cauchy--Binet-style determinant-of-product expansion
\lean{det\_mul\_rectangular}{SignedQMomentObstruction.lean:29} (definitions
\lean{truncatedMoment}{SignedQMomentObstruction.lean:108} and
\lean{hankelDet}{SignedQMomentObstruction.lean:117}). \ev{Lean} \scale{bounded}
\coord{p-adic}.

\textbf{(c) Precise scope — the point this section exists to make.} Despite
living in a file named alongside the corpus's ``obstruction'' modules, this
is explicitly \emph{not} phrased as a no-go in its own docstring: it is the
finite nonvanishing engine that a Hankel--Padé construction for \#249 would
\emph{consume}. The parity lemma and the Cauchy--Binet expansion use no
totient or Mersenne structure at all — pure linear algebra, fully
problem-agnostic, directly transplantable to any base-\(q\) dyadic-clearing
argument including \#257's series family. What is missing, and what
prevents this from closing anything, is a \emph{supply} theorem: no result
in the corpus proves that a unique-terminal configuration of this shape
exists cofinally for the totient-derived family. The obstruction, such as it
is, is not in this lemma but in the absence of its hypothesis's cofinal
supply — exactly the same shape of gap as \S\ref{sec:mahler-defect}'s Mahler
defect below.

\textbf{(d) Salvage.} The entire lemma is reusable machinery, not wreckage —
nothing here needs salvaging because nothing here failed. It stands ready
for whichever future construction can supply the missing cofinal
unique-terminal-configuration existence.

\subsection{The residual-gauge obstruction}

\textbf{(a) Route as conceived.} A ``first-harmonic pivot'' strategy for
\#249 (via exponential/character-sum cancellation,
\lean{FirstHarmonicPivot.lean}{FirstHarmonicPivot.lean}) might try to certify a nonzero determinant or
full-rank minor of a residual-weighted monomial matrix as sufficient
evidence that a genuine phase reconstruction (as opposed to a degenerate
locked configuration) has occurred.

\textbf{(b) Exact mechanism.}
\begin{prop}[Locked reconstruction preserves a nonzero minor]
For a monomial matrix \(\mathrm{residualMonomialMatrix}(e,r,z)(i,j) =
r(i,j)\cdot z(j)^{e(i)}\) with all residual entries \(z(j)\ne0\) and
\(e(i)=1\) for a distinguished row \(i\): \(\det(\mathrm{phasePowerMatrix})
\ne0 \implies \det(\mathrm{residualMonomialMatrix})\ne0\), \emph{and} the
distinguished row is identically \(1\) regardless.
\end{prop}
\lean{Erdos249257.locked\_reconstruction\_preserves\_nonzero\_minor}{ResidualGaugeObstruction.lean:94}
(row-dependent relaxation \lean{rowResidualMonomialMatrix\_reconstructs\_phasePowerMatrix}{ResidualGaugeObstruction.lean:124}; the
column-gauge identity underlying it,
\(\det(\mathrm{residualMonomialMatrix}) =
\det(\mathrm{phasePowerMatrix})\cdot\prod_j W(j)\), is
\lean{det\_residualMonomialMatrix}{ResidualGaugeObstruction.lean:47}; a row-dependent relaxation still fails at
\lean{rowReconstruction\_zero\_row\_one}{ResidualGaugeObstruction.lean:145}). \ev{Lean} \scale{uniform}
\coord{other:first-harmonic-phase}.

\textbf{(c) Precise scope.} Pure linear algebra over \(\mathbb C\), fully
generic in dimension \(d\) and exponent function \(e\), zero \#249-specific
content — this is a statement about \emph{residual-weighted monomial
matrices as a representation}, not about the totient object. It rules out
exactly one class of certificate: a residual-blind rank/determinant/
conditioning test, used alone, cannot distinguish genuine phase
reconstruction from an explicit ``locked'' degenerate configuration in which
a residual weighting scales every column by \(W(j)=z(j)^{-1}\), collapsing
the exponent-one row to the constant row \((1,\dots,1)\) while the
determinant stays nonzero whenever the base (Vandermonde-shaped)
\(\mathrm{phasePowerMatrix}\) determinant is nonzero. It does \emph{not}
rule out determinants carrying an additional arithmetic coupling identity
between rows.

\textbf{(d) Salvage.} A hard requirement for any future first-harmonic
determinant construction (Part E of the certificate bank,
\lean{irrational\_totient\_series\_of\_first\_harmonic\_norm\_gap}{FirstHarmonicPivot.lean:83}): the construction \emph{must}
couple rows by an extra arithmetic identity, not merely gauge-normalise
columns. The diagonal-gauge-action technique itself is reusable for any
determinant-based nonvanishing strategy in any coordinate.

\subsection{Cone flatness and cone nonflatness}

\textbf{(a) Route as conceived.} The diagonal certificate-supply obligation
(\lean{Erdos249257.irrational\_totient\_series\_iff\_lcm\_diagonal\_certificate\_supply}{LcmConeFlatness.lean:426}) concerns only the pair \((H,2H)\) at LCM
height \(H=\mathrm{periodLcm}(t)\). Two natural strengthenings: (i) show
rationality of \(S\) is even more constrained than a single diagonal pair —
it flattens an \emph{entire cone} of ratios; (ii) conversely, exhibit
\emph{nonflatness} on a wider ``menu'' of cone vertices as a cheaper
sufficient certificate than a single pairwise \(\mathrm{certifiedKill}\).

\textbf{(b) Exact mechanism — flatness (necessary consequence of
rationality).}
\begin{prop}[Wave-24 LCM-cone flatness]
\(\neg\mathrm{Irrational}(S) \implies \exists t_1,\ \forall t\ge t_1,\
\forall q,m:\mathbb N,\ 0<q \implies \mathrm{totientTail}(q\cdot
\mathrm{periodLcm}(t)+m\cdot\mathrm{periodLcm}(t)) -
\mathrm{totientTail}(q\cdot\mathrm{periodLcm}(t)) \in
\mathrm{range}((\uparrow):\mathbb Z\to\mathbb R)\).
\end{prop}
\lean{rational\_totient\_series\_forces\_lcm\_cone\_flatness}{CertificateKernel.lean:19002} (body
\texttt{LcmConeFlatness.lean}). \ev{Lean} \scale{uniform} (for all
sufficiently large \(t\), \emph{given} rationality) \coord{other:lcm-ray-totient}.
Rationality forces one fractional constant on the entire LCM cone
\(\{k\cdot\mathrm{periodLcm}(t):k\ge1\}\) at every scale \(t\ge t_1\), not
merely the single diagonal pair — the umbrella collapse theorem
\lean{irrational\_totient\_series\_of\_lcm\_cone\_window\_kill\_supply}{CertificateKernel.lean:19014} shows any
\(\mathrm{certifiedKill}\) \emph{anywhere} on the two-multiplier cone, at
arbitrarily large \(t\), forces \(\mathrm{Irrational}(S)\).

\textbf{(b\('\)) Exact mechanism — nonflatness (a sharper sufficient
certificate).}
\begin{prop}[Wave-25 cone-nonflat menu refuter]
For a nonempty menu \(Q\) of positive vertex multipliers with a
\emph{one-sided} floor \(\forall q\in Q,\ q\cdot H+L+2<2^L\) (half the
pairwise floor of \(\mathrm{certifiedKill}\)): if \(\mathrm{coneNonflatCert}\,
H\,L\,Q\) fires (an argmin/Helly-avoidance argument: the vertices cannot all
share one fractional part, else the minimal-deep-tail vertex would be a
common left endpoint of every arc), then \(\exists q_i,q_j\in Q,\
\mathrm{totientTail}(q_j H)-\mathrm{totientTail}(q_i H)\notin
\mathrm{range}((\uparrow):\mathbb Z\to\mathbb R)\).
\end{prop}
\lean{exists\_nonintegral\_pair\_of\_coneNonflatCert}{LcmConeNonflat.lean:126} (supply theorem
\lean{irrational\_totient\_series\_of\_lcm\_cone\_nonflat\_supply}{CertificateKernel.lean:19140}). \ev{Lean}
\scale{bounded} (proved for a given finite menu \(Q\); an
\emph{unbounded}-scale menu is the open sufficient target)
\coord{other:lcm-ray-totient}.

\textbf{(c) Precise scope.} Flatness is a genuine \emph{necessary
consequence} of rationality — it is a statement about what rationality of
\(S\) would force, conditional on rationality, not an unconditional fact
about \(S\). Nonflatness is an unconditional \emph{sufficient} certificate
mechanism, information-theoretically half the depth floor of the pairwise
\(\mathrm{certifiedKill}\) (one-sided versus two-sided radius charging), but
it has been proved only for individual finite menus \(Q\); no menu of
unbounded scale has been supplied.

Neither direction decides \#249: flatness
gives a contrapositive route (exhibit \emph{any} cone nonintegrality at
arbitrarily large \(t\) and rationality is refuted) and nonflatness sharpens
the certificate needed, but the underlying producer obligation — a
nonintegral tail difference at cofinally many scales — is unchanged.

Note
also the coordinate warning from the manuscript's own Appendix C
(\lean{rational\_totient\_series\_forces\_lcm\_cone\_flatness}{LcmConeFlatness.lean:90}): cone flatness is a
\emph{necessary} consequence, and only the sharper diagonal-pincer /
full-target-avoidance normal forms
(\lean{diagonal\_int\_iff\_foreignDiagonalDefect\_hits\_fullTarget}{DiagonalPincerDecomposition.lean:215},
\lean{irrational\_totientSeries\_of\_full\_target\_avoidance\_supply}{DiagonalPincerDecomposition.lean:290}) are an exact
\(\mathrm{iff}\) with irrationality.

\textbf{(d) Salvage.} The argmin/Helly-avoidance combinatorial-geometry
technique behind cone-nonflatness is a reusable pattern for any modular-residue
pincer argument with more than two points, in any coordinate. The rank-2
second-difference certificate family built on the same window apparatus was
independently tested and found \emph{not} to be a shortcut over rank 1
(measured, not merely conjectured: at \((h,N)=(1,8)\), rank-1 fires at depth
8 while no rank-2 certificate exists at depth \(\le8\);
\lean{totient\_tail\_rank\_two\_kill\_sound\_but\_not\_shallower\_cell}{CertificateKernel.lean:19090}) — an
explicitly flagged dead end not to re-attempt without new information.

\subsection{The Mahler defect: dyadic totient-kernel finite rank per level,
infinite rank overall}
\label{sec:mahler-defect}

\textbf{(a) Route as conceived.} A finite-linear-algebra shortcut to \#249:
find a bounded-depth linear relation among the dyadic totient-kernel
channels \(n\mapsto\varphi(2^jn+r)\) that would compress the infinite
family to a finite-dimensional space, from which a rationality-forcing
contradiction (or a genuine rank obstruction) might be extracted.

\textbf{(b) Exact mechanism.}
\begin{thm}[Infinite dyadic totient-kernel rank]
For every depth \(e\ge0\), the canonical family of \(2^e+1\) dyadic
totient-kernel channels (\(\mathtt{card\_totientCanonicalIndex}\),
\lean{card\_totientCanonicalIndex}{TotientMahlerDefect.lean:61}) is linearly independent over
\(\mathbb Q\), proved via a \texttt{SeparatedMinorCertificate} (an explicit
finite evaluation-point assignment with nonzero determinant, forced by CRT +
Dirichlet's theorem on primes in arithmetic progressions — one channel made
prime, every other channel forced through a fresh prime \(\equiv1\bmod\) a
large power of \(2\)).
\end{thm}
\lean{linearIndependent\_canonicalTotientKernelFamily}{TotientMahlerDefect.lean:935}; consequently the full
infinite family spans an infinite-dimensional \(\mathbb Q\)-space,
\lean{not\_finiteDimensional\_span\_fullTotientKernel}{TotientMahlerDefect.lean:1145}. A companion
impossibility result closes the natural repair attempt directly: a bounded
\emph{compressed-adjoint certificate} — a triple \((Q,A,\mathrm{boundary})\)
with \(Q\cdot v\cdot A=\mathrm{boundary}\), \(|\mathrm{boundary}|<Q\cdot v\),
\(A\ne0\) — is provably impossible,
\lean{false\_of\_compressedAdjointCertificate}{TotientMahlerDefect.lean:1183}. \ev{Lean}
\scale{uniform} (holds for every \(e\); existence side is unconditional, via
Mathlib's CRT + primes-in-AP machinery) \coord{binary-digit} (dyadic
totient-kernel — \emph{not} the Möbius coordinate).

\paragraph{Integral coordinates and the complete relation module.}
The independence statement also determines the arithmetic of the channel
lattice.  Fix $e\ge1$, write $F_{j,r}(n)=\varphi(2^j n+r)$, and let
$I_e=\{(j,r):0\le j\le e,\ 0\le r<2^j\}$.  Retain $F_{0,0}$,
$F_{1,0}$, and the odd-residue channels at levels $1,\ldots,e$; call
this retained index set $J_e$.  The totient identities reduce every omitted
channel to $F_i=a_iF_{j(i)}$ for a retained channel $j(i)$ and an integer
$a_i\ge0$.  Thus the retained family generates the lattice
$L_e=\sum_{i\in I_e}\mathbb Z F_i$.  Its rational independence gives
integer independence as well: every element of $L_e$ has unique integer
coordinates in these $2^e+1$ retained channels.

To describe \emph{all} integral relations, let $E_i$ denote the formal
coordinate vector in $\mathbb Z^{I_e}$ and use the evaluation map
\[\adjustbox{max width=\linewidth}{$\displaystyle
 \operatorname{ev}_e:\mathbb Z^{I_e}\longrightarrow L_e,
 \qquad (c_i)\longmapsto\sum_{i\in I_e}c_iF_i.
$}\]
For each omitted index put $R_i=E_i-a_iE_{j(i)}$.  Then
\[\adjustbox{max width=\linewidth}{$\displaystyle
 \ker(\operatorname{ev}_e)
   =\bigoplus_{i\in I_e\mathbin{\backslash} J_e}\mathbb Z R_i,
 \qquad
 \operatorname{rank}_{\mathbb Z}\ker(\operatorname{ev}_e)=2^e-2.
$}\]
Indeed, subtracting $\sum_{i\notin J_e}c_iR_i$ from a relation removes
all omitted coordinates.  The remainder is a relation among retained
channels, so independence makes it zero.  Conversely, the coefficient of
$E_i$ in $\sum_{h\notin J_e}b_hR_h$ is exactly $b_i$, proving uniqueness.
The decisive feature is the unit pivot $1$ in each omitted coordinate
(or $-1$ if a relation is oriented oppositely).  Elimination uses no division;
it therefore identifies the integral relation lattice, not merely its
rational span.  The depth-two example in the short note illustrates these
two-term rows.

The \href{https://github.com/wcook04/plectis-erdos/blob/25ef6245d15a47548c6926369ae8f1a0f0a14a80/ErdosProblems/Erdos249/PaperCompleteR8/KernelRelationBasis.lean#L398}{integral normal form} and the \href{https://github.com/wcook04/plectis-erdos/blob/25ef6245d15a47548c6926369ae8f1a0f0a14a80/ErdosProblems/Erdos249/PaperCompleteR8/FullKernelAssemblies.lean#L156}{full dyadic rational basis} are kernel-checked in Lean.
The formal statement gives the same construction for every integer base
$k\ge2$, with relation rank $\sum_{j=1}^{e-1}k^j$.  The full dyadic
assembly separately gives a basis for the infinite rational span and for
its finitely supported rational relations.  These finite and infinite
statements should not be confused with a claim about infinite sums of
relations.

\textbf{(c) Precise scope — coordinate-relative, stated explicitly by the
source module.} This does \emph{not} show irrationality of \(S\). It proves
the dyadic-kernel \emph{side} is infinite-rank; the module's own docstring
names the missing input as ``a rationality-side finite-rank compression, or
equivalent contradiction.''

The companion necessary-consequence-of-rationality
theorem, \(\neg\mathrm{Irrational}(S) \implies \exists v>0,\exists U:\mathbb
N\to\mathbb Z,\ \mathrm{IsTemperedBinaryOrbit}(\varphi,v,U)\wedge\forall
e,\ 2^e-1\le\mathrm{finrank}_{\mathbb Q}\,\mathrm{span}(\mathrm{range}(
\mathrm{canonicalCarryKernelFamily}(U,e)))\)
(\lean{not\_irrational\_totientSeries\_implies\_unbounded\_carryRank\_unconditional}{TotientCarryKernelRigidity.lean:300}),
shows the SCALE side is finished — the \(2^e-1\) floor holds for all \(e\)
with a uniform proof — but the opposite inequality (a rationality-side rank
\emph{upper} bound) is exactly what is missing, and its absence is the
entire open content here.

A \emph{third}, independent coordinate exhibits
the same coordinate-relative phenomenon: the Möbius-incidence companion
matrix \(U_N(i,j)=\mu((i+1)/(j+1))\) if \((j+1)\mid(i+1)\) else \(0\) is
lower triangular with diagonal \(1\), hence \(\det U_N=1\) for every \(N\),
hence its jet-evaluation map is injective — a finite ``incidence quotient''
compression is impossible at any finite horizon in the Möbius coordinate too
(\lean{incidenceMobius\_det\_eq\_one}{IncidenceQuotientHermitePade.lean:56}). Both natural finite
truncations — dyadic-residue and Möbius-incidence — turn out to have no
nontrivial kernel, in two completely different coordinates. That is strong
evidence any winning finite-linear-algebra shortcut, if one exists, must
live in a genuinely third coordinate or use an infinite/growing-parameter
construction, not evidence that no shortcut exists at all.

\textbf{(d) Salvage.} \lean{SeparatedMinorCertificate}{TotientMahlerDefect.lean:83} and
\lean{linearIndependent\_of\_separatedMinorCertificate}{TotientMahlerDefect.lean:91} are fully generic finite-rank
linear-independence infrastructure — the indexed family is a free variable,
reusable for any finite-rank obstruction argument on either problem (e.g.
testing independence of Mersenne-indexed sequences for \#257 via an explicit
nonzero minor). The CRT+Dirichlet evaluation-forcing technique (make one
channel prime, force every other channel through a fresh arithmetic-progression
prime) is a reusable construction pattern independent of the totient
specifics.

\subsection{Further closed routes, catalogued for completeness}

The following additional no-gos and route-pruning results were read in full
and are recorded here in compressed form; each follows the same (a)/(b)/(c)/(d)
discipline as above but is presented as a table row for space.

\begin{landscape}
\small
\begin{longtable}{@{}p{0.20\linewidth}p{0.38\linewidth}p{0.36\linewidth}@{}}
\toprule
\papertablehead
\textbf{Route} & \textbf{Mechanism \& site} & \textbf{Scope / salvage} \\
\midrule
\endfirsthead
\toprule
\papertablehead
\textbf{Route} & \textbf{Mechanism \& site} & \textbf{Scope / salvage} \\
\midrule
\endhead
\bottomrule
\endlastfoot
Full terminal dyadic staircase (annihilate every suffix letter mod growing
powers of 2) &
\(a\ge8\), room bound \(\Rightarrow\) \texttt{ActualLcmTerminalDyadicStaircase}
is false: the terminal letter would have to be positive (by short-window
positivity) and strictly below a wider-than-itself modulus while divisible
by it, forcing it to \(0\), contradicting positivity.
\lean{not\_actualLcmTerminalDyadicStaircase\_of\_room}{TotientActualLcmTopEdgeStaircase.lean:251}.
\ev{Lean} \scale{uniform} \coord{binary-digit} &
Kills the FULL terminal-staircase proof strategy outright and
unconditionally; the \emph{punctured} staircase (all but the last letter
vanish) survives and is pinned exactly to the half-turn value \(2^{m-1}\) at
the penultimate letter
(\lean{puncturedDyadicStaircase\_penultimate\_eq\_half}{TotientActualLcmTopEdgeStaircase.lean:1187}). The
generic mechanism (\(0<e<2^m\wedge2^m\mid e\Rightarrow e=0\)) is reusable
anywhere a positive quantity is asked to vanish mod a wider modulus. \\
LCM factor-ideal retention (keep only the homogeneous ``factor-only'' part
of an LCM-ray decomposition) &
An explicit nonzero all-horizon countermodel, built from multiples of
\(\varphi(\mathrm{periodLcm}(t))\), agrees with every exact whole-ray anchor
\(\Delta\varphi(H,qH)=\varphi(H)\) for \(2\le q<t\) and survives every finite
commensurate LCM-cube shift-polynomial, at every finite rank.
\lean{LcmFactorIdealPulseObstruction.lean:1--24}{LcmFactorIdealPulseObstruction.lean:1-24} (module docstring;
theorem bodies \ev{Cert}, DOCSTRING-SOURCED). \scale{uniform}
\coord{other:lcm-ray-totient} &
This is a \emph{representation}-level exclusion, not an object-level one: it
shows the homogeneous factor-ideal projection discards fresh Möbius-channel
information that can be adversarially reconstructed. Explicitly flagged
SYNTHETIC — no claim that the compensation letters occur as actual totient
differences. Salvage: a future factor-ideal argument must control the
\emph{fresh} channel via
\lean{totient\_eq\_sum\_mobiusTotientChannel}{LcmFactorIdealPulseObstruction.lean:328}, not just the old one. \\
Two-prime ``full-target diamond'' (four simultaneous ray hits certify more
than one) &
The four-hit diamond
\(\mathrm{Hit}(H)\wedge\mathrm{Hit}(pH)\wedge\mathrm{Hit}(qH)\wedge
\mathrm{Hit}(pqH)\) is logically equivalent to \(\mathrm{Hit}(H)\) alone, no
primality of \(p,q\) needed — hits transport multiplicatively for free
because \(\mathrm{Hit}(H)\iff\mathrm{IsIntegralValue}(\text{scaleDiagonalTailDifference}(H))\)
and integrality transports affinely along every ray \(H\mapsto kH\).
\lean{Erdos249257.fullTarget\_primeAdjunction\_diamond\_iff\_root}{FullTargetPrimeAdjunctionNoGo.lean:196}.
\ev{Lean} \scale{uniform} \coord{other:mobius-mersenne} &
The 4-condition diamond carries zero information beyond its base point — a
``curvature-zero collapse'' of a multi-point certificate. Any future
positive-holonomy argument needs an independently-defined projection of the
foreign state \emph{plus} control of the discarded complement; a bare
2/4-point diamond cannot supply new information beyond its own scale. \\
Fixed-precision local valuation-unit signature (2-adic ``tropical curvature
carry'' attack) &
For any finite word of odd-unit-part 2-adic symbols at fixed precision
\(u>0\) and any starting carry state \(e\), there exists a compatible carry
orbit realising the word with every intermediate state centred inside its
symbol's dyadic radius.
\lean{Erdos249257.TotientTailPeriodKiller.fixedPrecisionTropicalNoGo}{TropicalCurvatureCarry.lean:137}.
\ev{Lean} \scale{bounded} \coord{p-adic} &
Bounded LOCAL valuation-unit data at FIXED precision can never exclude all
finite centred carry completions — no obstruction is derivable from a
fixed-window local signature alone; growing precision, or extra arithmetic
coupling, is required. Pure 2-adic dynamics, zero totient content, reusable
against any similarly-shaped fixed-window carry-certificate attack. \\
Fixed-depth affine carry reset (distinguishing two carry histories after a
long common tail) &
For the affine binary orbit \(\mathrm{orbit}(0)=u_0\),
\(\mathrm{orbit}(n{+}1)=2\,\mathrm{orbit}(n)-a(n{+}1)\) with common forcing
word \(a\): \(\mathrm{orbit}_u(L)-\mathrm{orbit}_v(L)=2^L(u_0-v_0)\) exactly.
\lean{Erdos249257.affineBinaryOrbit\_mod\_twoPow\_eq}{GenericTailOrbitRigidity.lean:307}. \ev{Lean}
\scale{uniform} \coord{binary-digit} &
After \(L\) common steps the endpoint residue mod \(2^L\) is
\emph{independent} of the initial carry — long common tails erase
predecessor information exactly, not approximately. Directly relevant to any
argument hoping to distinguish two carry histories from a shared long
suffix; the terminal forcing word alone determines the residue. \\
Bounded finite-state / autonomous successor decoder for the binary carry &
For a balanced-pulse family at location \(m\) whose predecessor \texttt{State}
is constant across the whole family, no \(\mathrm{decode}:\mathrm{State}\to
\mathbb N\) recovers the radius parameter; any finite \texttt{Fintype State}
needs \(\mathrm{card}\ge\lfloor m/2\rfloor+2\), unbounded in \(m\).
\lean{Erdos249257.balancedPulse\_no\_autonomous\_decoder}{GenericTailOrbitRigidity.lean:247}. \ev{Lean}
\scale{uniform} \coord{binary-digit} &
Rules out, unconditionally, ANY proof strategy that defines a bounded or
autonomous ``carry state'' summarising pre-\(m\) history and claims it
exactly determines the post-\(m\) tail — independent of whether the
coefficients are \(\varphi\) or a Möbius-support indicator. A hard stop
against finite-automaton-computes-the-expansion style attacks on either
problem. \\
Period-4 sign weight as a route to A7's ``frequently nonzero'' hypothesis for
free &
The divisor coefficient of the period-4 sign weight
(\(w(a){=}1\) if \(a{\equiv}1\), \(-1\) if \(a{\equiv}3\), else \(0\), mod
4) vanishes identically at every \(n\equiv3\pmod4\), via the divisor-pairing
involution \(d\mapsto n/d\).
\lean{Erdos249257.intWeightedCoeff\_periodFourSignWeight\_eq\_zero\_of\_mod\_four\_eq\_three}{CertificateKernel.lean:14226}.
\ev{Lean} \scale{fixed} \coord{mobius-mersenne} &
Witnesses that Dirichlet-convolution coefficients CAN vanish identically on
an arithmetic progression, so the ``frequently nonzero'' hypothesis in the
nonnegative-periodic irrationality closure
(\lean{irrational\_intWeightedErdosSeries\_periodic\_of\_coeff\_nonneg\_of\_frequently\_ne\_zero}{CertificateKernel.lean:14583})
is not free — it is why the corpus lands a dichotomy for general periodic
signed weights rather than an unconditional supply theorem. \\
Certified finite-depth ``death'' certificates as a route to membership proofs
(both problems) &
\(\mathtt{CertifiedGreedyMersenneDeath}(x,\mathrm{level},\mathrm{lookahead})\)
is a decidable, finite-depth certificate proving non-membership only
(\lean{greedyMersenneDeath\_not\_mem}{GreedyAchievementSet.lean:1747}; witness
\lean{three\_fourths\_certifiedGreedyMersenneDeath}{GreedyAchievementSet.lean:1778}, \(3/4\) certified dead at
level 1). \ev{Cert} \scale{bounded} \coord{greedy-orbit} &
Structural one-sidedness, recurring across the WHOLE certified-kill family
(\lean{TotientTailPeriodKiller.lean}{TotientTailPeriodKiller.lean}, \lean{LcmConeFlatness.lean}{LcmConeFlatness.lean}, this one):
certified-death \(\Rightarrow\) non-membership, but absence of a found
certificate, or survival through any finite depth, proves \emph{nothing}
about membership or rationality. \\
Reconstructing one totient value from a two-tail absolute-value linear
combination &
For any finite \(w:\iota\to\mathbb Q\), \(x:\iota\to\mathbb N\) with
\(\sum w_i\varphi(x_i)=1\) (exact isolation), the crude two-tail cost
\(\sum|w_i|(2(x_i+1)+(x_i+2))\ge3\), using only \(\varphi(x)\le x\).
\lean{Erdos257PeriodNoncollapse.three\_le\_totientAdjugateTailCost}{TotientTailCarryPeriod.lean:266}
(closure \lean{not\_totientAdjugateTailCost\_lt\_one}{TotientTailCarryPeriod.lean:307}). \ev{Lean}
\scale{uniform} \coord{other:totient-adjugate-linear-algebra} &
The strategy needs cost \(<1\) and gets \(\ge3\) at ANY finite grid height —
kills the entire ``isolate one coefficient via absolute-value two-tail
bookkeeping'' family before it starts. The proof mechanism is problem-agnostic
(any \(c(n)\le n\) sequence gets the same floor by the triangle inequality)
and transfers verbatim to a \#257 reformulation. \\
Homogeneous Mersenne-primitive prime factor alone forces a contradiction &
For the completely-multiplicative control \(c(n)=n\) (zero totient content):
if \(K<q\), \(q\mid2^K-1\), the corresponding tail shift is integral at every
\(N\) yet \(q\nmid K\). \lean{idCoeff\_mersenneModulus\_does\_not\_annihilate\_tailShift}{TotientTailCarryPeriod.lean:342}.
\ev{Lean} \scale{n/a} \coord{other:mersenne-modulus} &
A primitive prime factor of the homogeneous multiplier \(2^K-1\) alone
supplies NO contradiction from tail integrality. A direct cross-problem
warning: \#257's denominators \(2^n-1\) are literally this \(q\mid2^K-1\)
object, so ``a large primitive Mersenne prime factor alone forces a
contradiction'' should be checked against this lemma before being attempted
on either problem. \\
Prime-power reduced-denominator unit-gap strengthening (rescuing extra Farey
lattice points) &
At a prime-power reduced denominator \(p^e\), the unit-gap refinement can
rescue at most ONE additional candidate lattice point beyond the ordinary
gap certificate. \lean{reducedDenominatorCandidateCount\_prime\_pow\_le\_one}{PrimitiveWeightCertificate.lean:111}
(iff-characterisation \lean{reducedDenominatorUnitGapCert\_prime\_pow\_iff}{PrimitiveWeightCertificate.lean:54}). \ev{Lean}
\scale{uniform} \coord{other:farey-gap} &
A formal ceiling on how much this specific refinement can ever buy — a
known dead end for extending the \(\approx7.96\times10^{34}\) Farey rung
(\S\ref{sec:adelic-height}) unboundedly. Do not re-attempt this exact
strengthening expecting more than +1 lattice point per prime power; the
underlying reduced-denominator-implies-unit-numerator technique is reusable
for other denominator-exclusion arguments. \\
Fixed rank-3 finite-difference kernel squeezing more than one dyadic bit of
curvature information &
\(\mathrm{fixedRankSecondDifference}(2^{n+1}\cdot2,\,2^{n+1}\cdot3)=-2^{n+1}\)
exactly, for all \(n\): the two-adic valuation gained by the primitive
kernel factor \(2\varphi(j)\) is exactly tight.
\lean{fixedRankSecondDifference\_dyadic\_fixture\_exact}{TotientFixedRankLcmAsymptotic.lean:462}. \ev{Lean}
\scale{n/a} \coord{binary-digit} &
Route-pruning ammunition: rules out extracting more than one dyadic bit from
the bare 3-rank affine kernel \((1,-2,1)\); any further progress on the
fixed-rank curvature route needs new arithmetic input, not a sharper
valuation squeeze from the same kernel. The affine rank-3 kernel
classification itself
(\lean{affine\_rank\_three\_kernel\_classification}{TotientFixedRankLcmAsymptotic.lean:338}, pure linear algebra
over \(\mathbb Z\), zero number-theoretic content) is fully reusable for any
3-point finite-difference argument on either problem. \\
\end{longtable}
\end{landscape}

\section{Why the bounds are load-bearing}
\label{sec:promotion-audit}

The near-miss interface between the corpus's proved results and the four
open supply obligations — \(\#249\)-supply, \(\#257\)-reset,
\(\#257\)-cofinal-rows, and \(\#257\)-universal — was swept for
\emph{promotion candidates}: proved theorems whose \emph{stated} hypotheses
looked, on a first read of the \texttt{yields} clause, close enough to the
open obligation's shape that a routine strengthening (raising a scale bound,
widening a hypothesis, or reindexing a quantifier) might close the gap
without new mathematics. Fifty-six near-miss rows were catalogued across the
four obligations; eighteen of them were flagged \texttt{promotable\_claimed:
True} on this first pass. \textbf{Every one of the eighteen proofs was then
opened and read in full}, not merely re-inspected at the statement level,
and the promotion claim was checked against the actual proof body. The
audit table below reports the verdict for all eighteen.

Of the eighteen, \textbf{sixteen are conclusively NOT\_PROMOTABLE}: reading
the proof body exposes a documented arithmetic, logical, or coordinate
obstruction that a routine strengthening cannot cross, and the row's own
record states exactly what new input would be required.

\textbf{Two}
(\texttt{NM-02}, \texttt{NM-03} in the \(\#257\)-universal obligation) are
genuine exceptions: the auditor found the promotion is available with, in
its own words, ``no new mathematics'' — a pure widening of an already-general
proof to a wider stated target. Crucially, \emph{neither of these two closes
anything}: they widen an equivalence or a corollary from one fixed target
value to a family of targets, without touching the missing arithmetic
content (a cofinal certificate supply, a rank upper bound, an anti-concentration
estimate) that every genuine closure of \(\#249\) or \(\#257\) requires.

So
while the promotion audit is not literally eighteen-for-eighteen at the level
of ``can this exact statement be re-derived for a wider scope,'' it is
eighteen-for-eighteen at the level that matters for this paper's obligations:
\textbf{no promotion candidate in the audited set supplies missing arithmetic
content toward closing \#249 or \#257}. The two mechanically-widenable rows
are recorded honestly below rather than folded into the sixteen, because
collapsing that distinction would itself be exactly the kind of
quantifier-slippage error this paper's accuracy rules forbid.

\begin{landscape}
\small
\begin{longtable}{@{}p{0.12\linewidth}p{0.18\linewidth}p{0.13\linewidth}p{0.51\linewidth}@{}}
\toprule
\papertablehead
\textbf{Obligation} & \textbf{Row / site} & \textbf{Gap kind} &
\textbf{Verdict and exact blocker} \\
\midrule
\endfirsthead
\toprule
\papertablehead
\textbf{Obligation} & \textbf{Row / site} & \textbf{Gap kind} &
\textbf{Verdict and exact blocker} \\
\midrule
\endhead
\bottomrule
\endlastfoot
249-supply & e1-companion,
\lean{exists\_certifiedKill\_of\_first\_harmonic\_gap}{FirstHarmonicGap.lean:128} & hypothesis strength &
Not promotable. The consumer already has the obligation's exact quantifier
shape; the missing input is a single unproved arithmetic fact — a
constant-saving first-harmonic (Weyl-sum) cancellation bound
\(\sum_{N\in[X,2X)}\cos(2\pi\cdot(\ldots)/2^L)\le(9/10)X\) — not proved at
any \(X,h,L\) anywhere in the corpus. \\
249-supply & SGN-01,
\lean{actualLcmTailDiff\_shift\_pos}{TotientActualLcmOrbitSign.lean:39} & hypothesis strength &
Not promotable. Positivity is exactly HALF of the needed certificate: the
companion theorem in the same file proves integrality forces the residue to
the TOP edge, not a kill — the lower half of the band is discharged
unconditionally and cofinally, the upper half is untouched. \\
249-supply & TE-04,
\lean{actualLcmTailDiff\_notMem\_int\_of\_topEdgeResidueGap}{TotientActualLcmTopEdgeStaircase.lean:1260} & hypothesis strength &
Not promotable. The consumer side is finished and cofinal (proved for
every \(a\ge8\)); the residual is purely the one-sided residue-gap
\emph{producer}, unconditionally unsupplied at even one large \(a\). \\
249-supply & TE-05-weakest,
\lean{PowerTwoFlexibleActualTerminalDominanceSupply}{TotientActualLcmTopEdgeStaircase.lean:1908} & hypothesis strength &
Not promotable. The exact identity pinning the target shows the dominance
inequality is equivalent to \(\mathrm{carryOrbit}\le0\), and the sign
machinery (SGN-02) already proves the true carry orbit is strictly
\emph{positive} under integrality — the corridor-escape branch the identity
naturally supplies is exactly the branch already eliminated. \\
249-supply & SEP-02,
\lean{abs\_actualLcmTailOrbit\_sub\_rawApprox\_lt}{TotientActualLcmOrbitSeparation.lean:141} & scale only &
Not promotable. The analytic tail is fully discharged with an explicit,
uniform error radius; only two exponents (\(a=4,6\)) have the resulting
finite distance-to-integer question verified unconditionally
(\lean{actualLcmTailOrbit\_four\_and\_six\_notMem\_int}{TotientActualLcmShortKill.lean:35}); nothing beyond
\(a=6\) is proved. \\
249-supply & b7,
\lean{upperHalfMersenneProduct\_lower\_bound}{MersenneShadowDenominatorGrowth.lean:60} & coordinate only &
Not promotable. The only proved unbounded-growth quantity in the
Möbius--Mersenne coordinate, uniform in \(t\), but there is no landed
transport from a denominator lower bound to \(\mathrm{certifiedKill}\) or to
totient-tail non-integrality — a different coordinate from the obligation. \\
249-supply & d4/d5,
\lean{linearIndependent\_canonicalTotientKernelFamily}{TotientMahlerDefect.lean:935}, \lean{not\_irrational\_totientSeries\_implies\_unbounded\_carryRank\_unconditional}{TotientCarryKernelRigidity.lean:300} &
coordinate only & Not promotable. The scale side of a second, independent
open obligation (unbounded carry-kernel rank forced by rationality) is
finished uniformly in \(e\); the opposite (rationality-side rank upper
bound) is missing, and one natural route to it is proved dead
(\lean{false\_of\_compressedAdjointCertificate}{TotientMahlerDefect.lean:1183}). \\
257-reset & rc-2 / F4-adjacent,
\lean{one\_third\_lt\_scaledMersenneTail\_floorWall}{HalfCylinderMiddleCarryLowerBound.lean:2308} & multiple &
Not promotable. The excess-side run-length law is only half-assembled: an
exact iterate identity exists, and an upper bound on the terminal excess
exists separately, but no theorem in Lean composes the two into an
excess-side run-length statement, even though every ingredient is present. \\
257-reset & E1,
\lean{SeamGreedyRemainderGapBound}{HalfCylinderSeamLimit.lean:192} & coordinate only &
Not promotable. Supplies an UPPER bound on the remainder (the ceiling half
of the picture); the obligation needs a LOWER bound on the deviation
magnitude — the two do not compose without an additional input. \\
257-reset & rc-8,
prose \emph{erdos257\_reset\_crossing\_unification\_2026\_07\_24.md}\,\S6 &
multiple & Not promotable. Two gaps at once: the sign law is
\ev{Cert}/empirical only (flagged \ev{Open} as an unconditional lemma, not
proved), and even if proved it delivers only the SIGN of the deviation, not
its magnitude, which is what the obligation needs. \\
257-cofinal-rows & TH-sharp-capacity-progress,
\lean{exists\_exactRowStrictUpperFill\_of\_skippedCoreSharpCapacity}{BooleanMobiusSkipRow.lean:238} & hypothesis strength &
Not promotable. Unconditional and strictly more general than every actual
consumer, but it needs a per-\(c\) input
(\(\mathrm{localBinarySuffix}\,D\,1\,(2c{-}2)<2^{c-2}\)) that the corpus
only ever supplies through a route which then rigidly collapses onto one
specific support family — the deficit hypothesis is a self-imposed
restriction of the consumers, and the general input itself is unsupplied. \\
257-cofinal-rows & precriticalSuffix\_lt\_of\_future\_skip\_after\_takenBlock,
\lean{halfGreedy\_precriticalSuffix\_lt\_of\_future\_skip\_after\_takenBlock}{BooleanMobiusCriticalCapacityCofinal.lean:603} &
hypothesis strength & Not promotable. Unconditional and uniform in both
parameters \((c,t)\); the sole missing input is an orbit-level skip-gap
bound on the rational half-greedy orbit, never proved and never previously
isolated as a target anywhere in the banks. \\
257-cofinal-rows & greedyHalfFrozenMargin\_nonneg,
\lean{exists\_greedyHalfFrozenMargin\_nonneg\_of\_excess\_neg}{HalfCylinderFiniteShadow.lean:1198} & scale only &
Not promotable. The margin-nonnegativity horizon is produced by a
\emph{non-effective} limit argument (an existential horizon from a
convergence fact); the exact-row route needs the SPECIFIC effective horizon
\(J=k-2\), and effectivising it is explicitly recorded as unfinished
bookkeeping, not new mathematics — but it is still unfinished. \\
257-cofinal-rows & C4a,
\lean{HalfCarryCofinalTerminalOnlyStrip}{TerminalOnlyCofinal.lean:34} & multiple &
Not promotable. Trades off against the obligation in the opposite
direction on \emph{support locality}: tolerates any support inside a depth
window but demands a carry inside a tight \(\sim2\sqrt M\) strip, where the
obligation tolerates an exponentially looser carry but demands the support
be confined to a lower window — the gap is support locality, not carry
size, and neither socket implies the other. \\
257-universal & NM-02,
\lean{infinite\_greedyMersenneSkippedSupport\_of\_half\_mem}{GreedyAchievementSet.lean:2547} & scale only &
\textbf{PROMOTABLE (no new mathematics), but does not close anything.}
The proof was opened and found already uniform in the target in every
internal step; only the statement's own quantifier is artificially pinned
to \(t=1/2\). Promoting it converts the whole half-greedy refutation machine
from a one-target equivalence to one covering every rational target — the
highest-value mechanical promotion in the bank — but it supplies no new
arithmetic content: it is a restatement, not a producer. \\
257-universal & NM-03,
\lean{finite\_boolSupport\_ne\_half}{HalfCarryReachability.lean:589} & scale only &
\textbf{PROMOTABLE (no new mathematics), but does not close anything.} The
headline is pinned to target \(1/2\) while its engine
(\lean{finiteErdosSum\_den\_odd}{DyadicPrefixCompression.lean:1513}) is already fully general over
finite supports; verbatim re-derivation at any even-denominator target is
available with zero new arithmetic. Paired with NM-02: universal \#257 at
\(b=2\) is false iff some rational \(t\) with even reduced denominator has
an infinitely-skipping greedy Mersenne orbit — a genuine reformulation, not
a resolution. \\
257-universal & NM-04,
\lean{dyadic\_support\_fraction\_reciprocalMass\_diverges\_or\_gt\_one}{RationalSupportCarrySkeleton.lean:2210} & coordinate only &
Not promotable (partial mechanical content only). The dyadic restriction
\(v=1\) is a call-site choice, not a proof constraint — the general-\(v\)
conclusion is already landed for arbitrary \((v,h,L,F)\) at
\lean{shiftedBooleanCollision\_wrap\_bound\_of\_common\_multiple}{RationalSupportCarrySkeleton.lean:1968} — but
the strict form needed for non-dyadic targets requires one genuinely new
auxiliary bound (a mean-least-residue estimate) not present in the corpus. \\
257-universal & NM-12,
\lean{half\_mem\_mersenneAchievementSet\_or\_exists\_fatal\_gap}{HalfCutLocator.lean:623} & scale only &
Not promotable (unverified, not merely blocked). The killer and the
dichotomy scaffold are already target-generic; the two endpoint kills that
pin the dichotomy to \(1/2\) generalise on their face, but their proof
bodies were explicitly \emph{not} read in this audit pass — recorded as a
strong promotability hypothesis, not a verified promotion, since a
\texttt{norm\_num}-style numeric step could hide a \(1/2\)-specific bound. \\
\end{longtable}
\end{landscape}

\subsection*{Classification of the sixteen genuine blockers}

Reading all eighteen proof bodies rather than trusting the \texttt{yields}
clause exposes a small, recurring taxonomy of blocker shapes, tagged by
\texttt{gap\_kind} in the table above:

\begin{itemize}[nosep]
\item \textbf{hypothesis\_strength} (6 rows: e1-companion, SGN-01, TE-04,
TE-05-weakest, TH-sharp-capacity-progress,
precriticalSuffix\_lt\_of\_future\_skip\_after\_takenBlock). The consumer
theorem is already at the obligation's exact quantifier shape; a single
named arithmetic fact — a Weyl-sum cancellation bound, a one-sided residue
gap, an orbit-level skip-gap bound — is missing and is not a re-derivation
of anything on disk.
\item \textbf{scale\_only} (5 rows: SEP-02, greedyHalfFrozenMargin\_nonneg,
NM-02, NM-03, NM-12). The mathematics is either already general or reduces
to a finite verified range; what is missing is either extending a finite
census (SEP-02) or effectivising a non-constructive existence bound
(greedyHalfFrozenMargin\_nonneg). NM-02 and NM-03 are the two genuine
exceptions where scale widening costs nothing — and, being restatements, buy
nothing toward closure either.
\item \textbf{coordinate\_only} (4 rows: b7, d4/d5, E1, NM-04). A quantity
is proved unconditionally in one coordinate (Möbius--Mersenne denominator
growth, dyadic-kernel rank, remainder upper bound, dyadic support-fraction
mass) with no landed transport into the coordinate the obligation is stated
in.
\item \textbf{multiple} (3 rows: rc-2/F4-adjacent, rc-8, C4a). More than
one of the above simultaneously — typically an unassembled composition of
two otherwise-landed halves, or two orthogonal deficiencies (empirical-only
plus sign-not-magnitude) stacked on the same row.
\end{itemize}

No row in the audited eighteen falls under \texttt{quantifier\_order} — every
row with a genuine quantifier-order mismatch (e.g. \texttt{a12}'s
\(\exists N\,\exists L\,\forall h\) versus the obligation's
\(\forall h\,\forall N_0\,\exists N\,\exists L\)) was judged NOT promotable
at first read and excluded from the eighteen entirely, rather than surviving
to a full proof-body audit. The eighteen audited here are precisely the rows
whose \emph{stated} hypotheses looked closest to the obligation; that even
this most-favourable subset yields sixteen genuine blockers and only two
costless-but-empty restatements is the paper's evidence that the corpus's
bounds are load-bearing in the strict sense: no amount of routine
strengthening of what is already proved, at the level the audit checked,
supplies the missing arithmetic content that \(\#249\) or \(\#257\)
actually needs.
